% Generated draft: requires mathematical and source review.
\documentclass[11pt]{article}
\usepackage[margin=25mm]{geometry}
\usepackage{fontspec}
\setmainfont{Latin Modern Roman}
\usepackage{amsmath,amssymb,mathtools,unicode-math}
\setmathfont{Latin Modern Math}
\usepackage{enumitem,needspace,xurl,hyperref}
\setlength{\parindent}{0pt}
\setlength{\parskip}{5pt}
\setlength{\emergencystretch}{3em}
\newcommand{\N}{\mathbb{N}}
\newcommand{\Z}{\mathbb{Z}}
\newcommand{\Q}{\mathbb{Q}}
\newcommand{\R}{\mathbb{R}}
\newcommand{\C}{\mathbb{C}}
\newcommand{\F}{\mathbb{F}}
\newcommand{\A}{\mathbb{A}}
\newcommand{\PP}{\mathbb{P}}
\newcommand{\E}{\mathbb{E}}
\newcommand{\eps}{\varepsilon}
\newcommand{\dd}{\,\mathrm{d}}
\begin{document}
\section*{UnsolvedMath ID: 30005528}
\textbf{Source problem code:} catalogue entry
\subsection{Definitions and mathematical statement}
The supplied reconstructed formulation concerns a linear domino game \(G=(\Sigma,T,s)\), where \(\Sigma\) is an infinite compact topological symbol space and \(s\) is a continuous size function. The record does not define \(T\), admissible domino sequences, the domain and codomain of \(s\), or the convention for average size. Consequently, the following is a conditional formulation, pending recovery of those definitions from the cited report.

Write \(X_G\) for the intended set of admissible sequences and \(A_G(x)\) for the intended average size of \(x\), assuming that this average is defined and comparable for all \(x\in X_G\). Define
\[
M_G=\{x\in X_G: A_G(x)\leq A_G(y)\text{ for every }y\in X_G\}.
\]
A sequence \(x=(x_i)\) is periodic if there is an integer \(p\geq1\) such that \(x_{i+p}=x_i\) for every applicable index \(i\). The source does not specify whether sequences are one-sided or two-sided. Let \(P_G\subseteq X_G\) denote the periodic admissible sequences.

The question suggested by the reconstructed statement and its proposed counterexample is whether
\[
M_G\cap P_G\neq\varnothing
\qquad\text{for every game }G\text{ satisfying the intended hypotheses}.
\]
If existence of minimizers is an intended standing assumption, a counterexample would satisfy \(M_G\neq\varnothing\) and \(M_G\cap P_G=\varnothing\). The supplied hypotheses do not establish that standing assumption. Moreover, requiring every minimizer to be periodic would instead mean \(M_G\subseteq P_G\), a different assertion.

There is an unresolved conflict in the optimization target: the reconstructed statement asks about minimizers, whereas the original excerpt asks, ``Can a maximizer be periodic?'' The supplied literature assessment also refers to maximizers. The intended extremum and the scope of the periodicity quantifier therefore require confirmation; the displayed formulation records only the minimization interpretation suggested by the reconstruction.
\subsection{Short English statement}
The reconstructed question asks whether every intended compact linear domino game with infinitely many symbols has a periodic sequence attaining minimum average size. The source lacks the game and averaging definitions and conflicts over minimization versus maximization, so this formulation remains conditional and requires human mathematical review.
\subsection{Sources}
\begin{itemize}
\item Supplied UnsolvedMath source record in the local registry, permanent ID \(30005528\), source problem code \texttt{catalogue entry}, title \emph{Periodic Minimizers in Compact Linear Domino Games}.
\item \emph{Group Actions and Harmonic Analysis in Number Theory} (2023), Oberwolfach Reports, pp.\ 1195--1262, DOI \texttt{10.4171/owr/2023/22}. Supplied source URL and literature-source link, labeled ``Oberwolfach problem report'': \url{https://doi.org/10.4171/owr/2023/22}.
\item The supplied excerpt credits Jens Marklof as chair and Paul Nelson as note taker. No proposer is identified for this problem. Its trailing text begins a separate problem attributed to Javier Fres\'an and does not supply definitions for the domino question.
\item The registry labels the problem open and the reconstructed statement \texttt{corrected\_verified}; it reports that the reconstruction was checked against the official report. These are supplied curation claims, not independent verification here.
\item Literature assessment dated exactly \texttt{2026-08-22}: the supplied curation describes the periodicity question for a maximizer with infinite \(\Sigma\) as open and reports finding no general existence or nonexistence theorem. This is supplied dated curation, not an independently verified literature assessment; the record expressly requests independent verification of time-sensitive status and cited claims.
\end{itemize}
\end{document}
